\begin{notechapter}
\chapter{Optimization Algorithms for Training}\label{ch:optimizers}
\chaptersubtitle{Gradient descent, momentum, adaptive methods and learning-rate schedules}

\begin{chapterintro}
Backpropagation gives the gradient; an optimizer decides what to do with it. Plain gradient
descent is slow when the curvature of the loss differs a lot between directions, and noisy
when the gradient is estimated from mini-batches. Momentum accelerates it, adaptive methods
rescale each coordinate, and schedules vary the step size over time. This chapter analyses
them on the quadratic model, where everything can be computed, and compares them on one
example. You will learn
\begin{itemize}[nosep]
  \item how the condition number limits gradient descent, and what the Hessian says (and
        does not say) about a critical point;
  \item heavy-ball and Nesterov momentum, and why momentum improves the rate from
        $\frac{\kappa - 1}{\kappa + 1}$ to $\frac{\sqrt\kappa - 1}{\sqrt\kappa + 1}$;
  \item AdaGrad, RMSProp and Adam, the bias correction of Adam, and AdamW;
  \item learning-rate schedules: step, exponential, cosine, warm-up, one-cycle;
  \item quasi-Newton methods (BFGS, L-BFGS) and when they are used.
\end{itemize}
\prerequisites{\cref{ch:feedforward-networks}; eigenvalues of symmetric matrices; the
geometry of the gradient and the zig-zag example (\cref{ch:gradient-tangent-space}).}
\end{chapterintro}

\section{Gradient descent and curvature}\label{optimizers:sec:gd}

\paragraph{The Hessian.} Near a point $\theta$, a twice differentiable loss is approximated by its
second-order Taylor expansion
\[
  \ell(\theta + v) = \ell(\theta) + \nabla\ell(\theta)^\top v + \tfrac12v^\top H(\theta)\,v + o(\norm v^2),\qquad H_{jk} = \frac{\partial^2\ell}{\partial\theta_j\partial\theta_k} .
\]
$H$ is symmetric; its eigenvalues are the curvatures along its eigenvectors. At a
\emph{critical point} ($\nabla\ell(\theta) = 0$): if $H$ is positive definite, $\theta$ is a strict local minimum; if
negative definite, a strict local maximum; if $H$ has eigenvalues of both signs, a saddle
point; if $H$ is only semidefinite, the second order does not decide. Away from a critical
point the Hessian says nothing of this kind: the gradient term dominates. In large
networks, critical points with many negative eigenvalues (saddles) are far more common than
local maxima, and the Hessian at the end of training typically has a few large
eigenvalues and a bulk near zero.

\paragraph{The quadratic model.} Take $\ell(\theta) = \frac12(\theta - \theta^*)^\top H(\theta - \theta^*)$ with $H$ symmetric, eigenvalues
$0<\mu = \lambda_{\min}\leq\dots\leq\lambda_{\max} = L$, and \emph{condition number} $\kappa = L/\mu$. Gradient descent
$\theta_{k+1} = \theta_k - \eta H(\theta_k - \theta^*)$ gives, for the error $e_k = \theta_k - \theta^*$,
\[
  e_{k+1} = (I - \eta H)\,e_k,\qquad e_k = (I - \eta H)^ke_0 .
\]
In the eigenbasis of $H$, each coordinate is multiplied by $1 - \eta\lambda_i$ at every step.

\begin{theorem}[Gradient descent on a quadratic]\label{optimizers:thm:gd}
Gradient descent converges for every $e_0$ if and only if $0<\eta<2/L$, and
$\norm{e_k}\leq\rho(\eta)^k\norm{e_0}$ with $\rho(\eta) = \max_i\abs{1 - \eta\lambda_i}$. The rate is best for $\eta = 2/(L + \mu)$, where
\[
  \rho = \frac{L - \mu}{L + \mu} = \frac{\kappa - 1}{\kappa + 1} .
\]
\end{theorem}

\begin{proof}
$I - \eta H$ is symmetric with eigenvalues $1 - \eta\lambda_i$, so its norm is $\rho(\eta)$, and $e_k\to0$ for all $e_0$ iff
all $\abs{1 - \eta\lambda_i}<1$, i.e.\ $0<\eta\lambda_i<2$ for all $i$. The function $\eta\mapsto\max(\abs{1 - \eta\mu},\abs{1 - \eta L})$ is minimal where the two
terms are equal with opposite signs, $1 - \eta\mu = \eta L - 1$, i.e.\ $\eta = 2/(L + \mu)$, where it equals $(L - \mu)/(L + \mu)$.
\end{proof}

The step size is limited by the \emph{largest} curvature (beyond $2/L$ the steep direction
diverges), and the speed by the \emph{smallest}: along the flat direction the error is
multiplied by $1 - \eta\mu\approx1 - 2/\kappa$ per step, so about $\kappa/2$ steps divide it by $e$. When $\kappa$ is large the
iterates oscillate across the steep direction while crawling along the flat one: the
zig-zag of \cref{gradient-tangent-space:sec:zigzag}. Newton's method replaces $\eta$ by $H^{-1}$ and
converges in one step on a quadratic; it is gradient descent for the metric defined by $H$
(\cref{gradient-tangent-space:sec:metric}).

\begin{example}[A running example]\label{optimizers:ex:quadratic}
$\ell(\theta) = \frac12(\theta_1^2 + 25\theta_2^2)$, so $\mu = 1$, $L = 25$, $\kappa = 25$, from $\theta_0 = (10, 1)$, stopping when the loss is below $10^{-6}$
times its initial value $62.5$. With $\eta = 1/L$, $\theta_2$ is zeroed in one step, but $\theta_1$ shrinks by $0.96$ per
step: $167$ steps. With the optimal $\eta = 2/26$, the rate is $\frac{24}{26} = 0.923$ and $87$ steps are needed; the
second coordinate changes sign at every step (\cref{optimizers:fig:paths}).
\end{example}

\paragraph{Stochastic gradients.} With a mini-batch gradient $g = \nabla\ell + \xi$, where $\xi$ has mean $0$ and
covariance $\Sigma/B$ (\cref{feedforward-networks:prop:minibatch}), the iterates do not converge to $\theta^*$ with a
constant step: they fluctuate in a region whose size grows with $\eta$ and with the noise. Robbins
and Monro (\emph{A Stochastic Approximation Method}, Annals of Mathematical Statistics,
1951) showed that steps with $\sum_k\eta_k = \infty$ and $\sum_k\eta_k^2<\infty$, such as $\eta_k\propto1/k$, give convergence under suitable
assumptions; in practice, schedules decrease $\eta$ in stages (\cref{optimizers:sec:schedules}).

\section{Momentum}\label{optimizers:sec:momentum}

\paragraph{Heavy ball.} Polyak's method (\emph{Some Methods of Speeding up the Convergence of
Iteration Methods}, USSR Computational Mathematics and Mathematical Physics, 1964) adds a
fraction of the previous displacement:
\[
  \theta_{k+1} = \theta_k - \eta\nabla\ell(\theta_k) + \beta(\theta_k - \theta_{k-1}),\qquad0\leq\beta<1 .
\]
With a velocity $v$, two equivalent forms are common:
\begin{align*}
  &\text{(PyTorch)} & v_{k+1} &= \beta v_k + \nabla\ell(\theta_k), & \theta_{k+1} &= \theta_k - \eta v_{k+1};\\
  &\text{(classical)} & v_{k+1} &= \beta v_k - \eta\nabla\ell(\theta_k), & \theta_{k+1} &= \theta_k + v_{k+1} .
\end{align*}
They give the same iterates for a constant $\eta$, and differ when $\eta$ changes over time. The
velocity is an exponentially weighted sum of past gradients, $v_{k+1} = \sum_{i\leq k}\beta^{k-i}\nabla\ell(\theta_i)$: components
that keep their sign (the flat direction) accumulate up to a factor $1/(1 - \beta)$, while components
that alternate (the oscillation across a valley) cancel.

\begin{theorem}[Heavy ball on a quadratic]\label{optimizers:thm:hb}
On the quadratic model, with
\[
  \eta = \frac4{(\sqrt L + \sqrt\mu)^2},\qquad\beta = \Bigl(\frac{\sqrt\kappa - 1}{\sqrt\kappa + 1}\Bigr)^2,
\]
the error of the heavy-ball method decreases like $k\,r^k$ with $r = \sqrt\beta = \dfrac{\sqrt\kappa - 1}{\sqrt\kappa + 1}$.
\end{theorem}

\begin{proof}
In the eigenbasis, each coordinate follows $x_{k+1} = (1 + \beta - \eta\lambda)x_k - \beta x_{k-1}$, a linear recurrence with
characteristic polynomial $r^2 - (1 + \beta - \eta\lambda)r + \beta$. The product of its roots is $\beta$. If the discriminant is
$\leq0$, the roots are complex conjugate (or double), both of modulus $\sqrt\beta$, and $\abs{x_k}$ is at most a
constant times $k\beta^{k/2}$. The discriminant $(1 + \beta - \eta\lambda)^2 - 4\beta$ is $\leq0$ iff $(1 - \sqrt\beta)^2\leq\eta\lambda\leq(1 + \sqrt\beta)^2$. Requiring
this for all $\lambda\in[\mu,L]$ with the smallest $\beta$ gives $\eta\mu = (1 - \sqrt\beta)^2$ and $\eta L = (1 + \sqrt\beta)^2$. Dividing,
$\sqrt\kappa = (1 + \sqrt\beta)/(1 - \sqrt\beta)$, so $\sqrt\beta = (\sqrt\kappa - 1)/(\sqrt\kappa + 1)$; adding the square roots, $\sqrt\eta(\sqrt L + \sqrt\mu) = 2$.
\end{proof}

The number of steps to divide the error by $e$ goes from about $\kappa/2$ to about $\sqrt\kappa/2$: for $\kappa = 10^4$,
from $5000$ to $50$. In \cref{optimizers:ex:quadratic}, $r = \frac46 = 0.667$, and $25$ steps suffice
(the factor $k$ in $kr^k$ explains why it is more than the $17$ that $r^k$ alone would suggest). The
guarantee is specific to quadratics: on general smooth strongly convex functions the
heavy-ball method with these parameters can fail to converge (Lessard, Recht and Packard,
\emph{Analysis and Design of Optimization Algorithms via Integral Quadratic Constraints},
SIAM Journal on Optimization, 2016).

\paragraph{Nesterov momentum.} Nesterov (1983) evaluates the gradient at the point the
momentum is about to reach:
\[
  v_{k+1} = \beta v_k + \nabla\ell(\theta_k - \eta\beta v_k),\qquad\theta_{k+1} = \theta_k - \eta v_{k+1} .
\]
With $\eta = 1/L$ and $\beta = (\sqrt\kappa - 1)/(\sqrt\kappa + 1)$, the loss gap decreases like $(1 - 1/\sqrt\kappa)^k$ for every $L$-smooth,
$\mu$-strongly convex function, not only quadratics (Nesterov, \emph{Introductory Lectures on
Convex Optimization}, 2004, §2.2). On the running example it needs $41$ steps. In deep
learning both forms are used with $\beta\approx0.9$; their benefit there, with noisy gradients and a
non-convex loss, is empirical (Sutskever, Martens, Dahl and Hinton, \emph{On the Importance of
Initialization and Momentum in Deep Learning}, ICML, 2013). PyTorch implements
\texttt{SGD(..., momentum=0.9, nesterov=True)} in an equivalent form that evaluates
gradients at the current iterate.

\section{Adaptive learning rates}\label{optimizers:sec:adaptive}

Adaptive methods give each coordinate its own step size, based on the magnitudes of its
past gradients. Below, operations on vectors are coordinatewise, $g_k = \nabla\ell(\theta_k)$ (or its
mini-batch estimate), and $\epsilon\approx10^{-8}$ prevents division by zero.

\paragraph{AdaGrad} (Duchi, Hazan and Singer, JMLR, 2011) divides by the root of the
accumulated squares:
\[
  G_k = G_{k-1} + g_k^2,\qquad\theta_{k+1} = \theta_k - \frac\eta{\sqrt{G_k} + \epsilon}\,g_k .
\]
Coordinates with rare or small gradients get large steps, which suits sparse features. But
$G_k$ only grows, so the steps eventually become too small for non-convex training.

\paragraph{RMSProp} (Tieleman and Hinton, lecture 6.5 of the Coursera course \emph{Neural
Networks for Machine Learning}, 2012) replaces the sum by an exponential moving average,
$G_k = \gamma G_{k-1} + (1 - \gamma)g_k^2$ with $\gamma\approx0.9$, so the scale adapts to recent gradients only.

\paragraph{Adam} (Kingma and Ba, \emph{Adam: A Method for Stochastic Optimization}, ICLR,
2015) combines a moving average of gradients (momentum) with RMSProp's scaling:
\begin{align*}
  m_k &= \beta_1m_{k-1} + (1 - \beta_1)g_k, & v_k &= \beta_2v_{k-1} + (1 - \beta_2)g_k^2, & m_0 = v_0 = 0,\\
  \hat m_k &= \frac{m_k}{1 - \beta_1^k}, & \hat v_k &= \frac{v_k}{1 - \beta_2^k}, & \theta_{k+1} = \theta_k - \eta\frac{\hat m_k}{\sqrt{\hat v_k} + \epsilon},
\end{align*}
with defaults $\beta_1 = 0.9$, $\beta_2 = 0.999$, $\eta = 10^{-3}$.

\begin{proposition}[Bias correction]\label{optimizers:prop:bias}
If the $g_k$ have a common mean $\bar g$, then $\E[m_k] = (1 - \beta_1^k)\bar g$, so $\E[\hat m_k] = \bar g$; the same holds for $v_k$
and the second moment.
\end{proposition}

\begin{proof}
Unrolling, $m_k = (1 - \beta_1)\sum_{i=1}^k\beta_1^{k-i}g_i$, and $(1 - \beta_1)\sum_{i=1}^k\beta_1^{k-i} = 1 - \beta_1^k$.
\end{proof}

Without the correction, the averages started at $0$ are biased towards $0$ during the first
$1/(1 - \beta)$ steps, which for $\beta_2 = 0.999$ is a thousand steps.

\paragraph{What Adam does.} At the first step $\hat m_1 = g_1$ and $\hat v_1 = g_1^2$, so the update is $-\eta g_1/(\abs{g_1} + \epsilon)\approx-\eta\sign(g_1)$:
every coordinate moves by $\eta$, whatever the size of its gradient. Later, the ratio $\hat m/\sqrt{\hat v}$ is
close to $\pm1$ when a coordinate's gradient is consistent and small when it fluctuates. The
update is invariant to multiplying the gradient by a fixed positive diagonal matrix, so
badly scaled coordinates are not a problem in themselves. On the running example with
$\eta = 0.3$, Adam moves $\theta_1$ by about $0.3$ per step at first, regardless of curvature, and needs $122$
steps; its step size, not the curvature, sets its speed (\cref{optimizers:fig:paths}).

\begin{figure}[ht]
\centering
\begin{tikzpicture}
\begin{axis}[width=14cm, height=6cm, xmin=-1.5, xmax=10.8, ymin=-2.1, ymax=2.1,
  xlabel={$\theta_1$}, ylabel={$\theta_2$}, axis on top,
  legend style={font=\small, draw=none, at={(0.5,1.03)}, anchor=south},
  legend cell align=left, legend columns=3]
  \foreach \c in {0.5,2,8,25,62.5}{
    \addplot[gray!40, domain=0:360, samples=90, forget plot]
      ({sqrt(2*\c)*cos(x)}, {sqrt(2*\c/25)*sin(x)});
  }
  \addplot[bookaccent, thick, mark=*, mark size=1pt] coordinates {(10.000,1.000) (9.231,-0.923) (8.521,0.852) (7.865,-0.787) (7.260,0.726) (6.702,-0.670) (6.186,0.619) (5.710,-0.571) (5.271,0.527) (4.866,-0.487) (4.491,0.449) (4.146,-0.415) (3.827,0.383) (3.533,-0.353) (3.261,0.326) (3.010,-0.301) (2.778,0.278) (2.565,-0.256) (2.367,0.237) (2.185,-0.219) (2.017,0.202) (1.862,-0.186) (1.719,0.172) (1.587,-0.159) (1.465,0.146) (1.352,-0.135)};
  \addlegendentry{gradient descent}
  \addplot[bookpitfall, thick, mark=square*, mark size=1pt] coordinates {(10.000,1.000) (8.889,-1.778) (7.407,1.926) (5.926,-1.778) (4.609,1.514) (3.512,-1.229) (2.634,0.966) (1.951,-0.741) (1.431,0.559) (1.040,-0.416) (0.751,0.306) (0.540,-0.224) (0.385,0.162) (0.274,-0.116) (0.194,0.083) (0.137,-0.059) (0.096,0.042) (0.068,-0.030) (0.047,0.021) (0.033,-0.015) (0.023,0.010) (0.016,-0.007) (0.011,0.005) (0.008,-0.004) (0.005,0.002) (0.004,-0.002)};
  \addlegendentry{heavy ball}
  \addplot[bookdef, thick, mark=triangle*, mark size=1.2pt] coordinates {(10.000,1.000) (9.700,0.700) (9.400,0.407) (9.101,0.132) (8.802,-0.110) (8.505,-0.300) (8.208,-0.424) (7.913,-0.481) (7.619,-0.481) (7.327,-0.434) (7.037,-0.355) (6.750,-0.255) (6.465,-0.144) (6.182,-0.032) (5.903,0.073) (5.627,0.161) (5.355,0.226) (5.087,0.265) (4.823,0.277) (4.563,0.263) (4.309,0.228) (4.059,0.177) (3.814,0.115) (3.575,0.049) (3.342,-0.016) (3.114,-0.074)};
  \addlegendentry{Adam}
\end{axis}
\end{tikzpicture}
\caption{The first $25$ steps on $\ell = \frac12(\theta_1^2 + 25\theta_2^2)$ from $(10,1)$, with level sets of $\ell$. Gradient
descent ($\eta = 2/26$) zig-zags across the narrow direction and advances slowly; the heavy ball
($\eta = 4/36$, $\beta = 4/9$) overshoots but covers the flat direction fast; Adam ($\eta = 0.3$) moves at a nearly
constant speed of about $\eta$ per coordinate.}
\label{optimizers:fig:paths}
\end{figure}

\begin{center}
\begin{tabular}{lll}
\toprule
method & parameters & steps to $10^{-6}$ \\
\midrule
gradient descent & $\eta = 1/L$ & $167$ \\
gradient descent & $\eta = 2/(L + \mu)$ & $87$ \\
heavy ball & \cref{optimizers:thm:hb} & $25$ \\
Nesterov & $\eta = 1/L$, $\beta = (\sqrt\kappa - 1)/(\sqrt\kappa + 1)$ & $41$ \\
Adam & $\eta = 0.1$, $0.3$, $1$ & $266$, $122$, $111$ \\
\bottomrule
\end{tabular}
\end{center}

On this deterministic, well-understood problem, tuned momentum wins. Adam's strength is
elsewhere: robustness to badly scaled parameters and to the choice of $\eta$ on large noisy
problems, which is why it is the default for transformers.

\paragraph{Caveats and AdamW.} Adam can fail to converge even on simple convex problems,
because $\hat v$ can decrease and enlarge steps (Reddi, Kale and Kumar, \emph{On the Convergence of
Adam and Beyond}, ICLR, 2018, who propose AMSGrad with a non-decreasing $v$). Adding an $L^2$
penalty $\frac\lambda2\norm\theta^2$ to the loss puts $\lambda\theta$ inside the gradient, where it is divided by $\sqrt{\hat v}$: coordinates with
large gradients are barely regularized. AdamW (Loshchilov and Hutter, \emph{Decoupled Weight
Decay Regularization}, ICLR, 2019) applies the decay outside the adaptive step,
\[
  \theta_{k+1} = \theta_k - \eta\Bigl(\frac{\hat m_k}{\sqrt{\hat v_k} + \epsilon} + \lambda\theta_k\Bigr),
\]
and generalizes better in practice. For plain SGD the two forms coincide. Weight decay
shrinks towards the origin of parameter space, a choice discussed in
\cref{gradient-tangent-space:ex:decay}.

\section{Learning-rate schedules}\label{optimizers:sec:schedules}

A schedule $\eta_t$, with $t$ the step or the epoch, combines a large rate early (fast progress,
and noise that explores) with a small rate late (to settle below the noise floor). Common
choices, with initial rate $\eta_0$ and $T$ total steps:
\begin{itemize}[nosep]
  \item \textbf{step decay} $\eta_t = \eta_0\gamma^{\lfloor t/s\rfloor}$, e.g.\ divide by $10$ every $s$ epochs;
  \item \textbf{exponential decay} $\eta_t = \eta_0\gamma^t$;
  \item \textbf{cosine} $\eta_t = \eta_{\min} + \frac12(\eta_0 - \eta_{\min})\bigl(1 + \cos(\pi t/T)\bigr)$, from $\eta_0$ to $\eta_{\min}$ (Loshchilov and Hutter,
        \emph{SGDR: Stochastic Gradient Descent with Warm Restarts}, ICLR, 2017, who also restart
        it periodically);
  \item \textbf{linear warm-up} $\eta_t = \eta_0t/T_w$ for $t\leq T_w$, then another schedule: large rates are unstable in the
        first steps, when the weights and Adam's statistics are far from typical values (Goyal
        et al., \emph{Accurate, Large Minibatch SGD: Training ImageNet in 1 Hour}, 2017);
  \item \textbf{one-cycle}: increase from a small rate to a large one, then decrease below the
        start, over a single cycle (Smith and Topin, \emph{Super-Convergence}, 2019).
\end{itemize}

\begin{figure}[ht]
\centering
\begin{tikzpicture}
\begin{axis}[width=10cm, height=5cm, xmin=0, xmax=100, ymin=0, ymax=1.1, xlabel={step $t$ (\% of training)},
  ylabel={$\eta_t/\eta_0$}, legend style={font=\small, draw=none, at={(1.02,0.5)}, anchor=west}, legend cell align=left]
  \addplot[very thick, bookaccent, const plot] coordinates {(0,1) (30,0.1) (60,0.01) (90,0.001) (100,0.001)};
  \addlegendentry{step ($\gamma = 0.1$, $s = 30$)}
  \addplot[very thick, bookexample, dashed, domain=0:100, samples=101] {0.5*(1+cos(180*x/100))};
  \addlegendentry{cosine}
  \addplot[very thick, bookdef, domain=0:100, samples=201] {x<10 ? x/10 : 0.5*(1+cos(180*(x-10)/90))};
  \addlegendentry{warm-up + cosine}
\end{axis}
\end{tikzpicture}
\caption{Three schedules over a training run.}
\label{optimizers:fig:schedules}
\end{figure}

In PyTorch a scheduler wraps the optimizer and is stepped after each epoch or batch:
\begin{lstlisting}[language=Python]
opt = torch.optim.AdamW(model.parameters(), lr=3e-4, weight_decay=0.01)
sched = torch.optim.lr_scheduler.CosineAnnealingLR(opt, T_max=num_steps)
for xb, yb in batches:          # num_steps iterations
    opt.zero_grad()
    loss_fn(model(xb), yb).backward()
    opt.step()
    sched.step()                # update the learning rate
\end{lstlisting}

\section{Second-order and quasi-Newton methods}\label{optimizers:sec:quasi-newton}

Newton's step $\theta - H^{-1}\nabla\ell$ needs the Hessian, which has $p^2$ entries for $p$ parameters: impossible
to store for a network, and not positive definite away from a minimum. \emph{Quasi-Newton}
methods build an approximation $B_k\approx H$ from gradients alone. With $s_k = \theta_{k+1} - \theta_k$ and
$y_k = \nabla\ell(\theta_{k+1}) - \nabla\ell(\theta_k)$, they impose the \emph{secant equation} $B_{k+1}s_k = y_k$, which a quadratic's Hessian
satisfies exactly. The BFGS update (Broyden, Fletcher, Goldfarb and Shanno, 1970) is
\[
  B_{k+1} = B_k + \frac{y_ky_k^\top}{y_k^\top s_k} - \frac{B_ks_ks_k^\top B_k}{s_k^\top B_ks_k},
\]
and the step is $-\alpha_kB_k^{-1}\nabla\ell(\theta_k)$ with a line search. $B_{k+1}$ stays positive definite as long as the
\emph{curvature condition} $y_k^\top s_k>0$ holds, which a line search satisfying the Wolfe conditions
guarantees (Nocedal and Wright, \emph{Numerical Optimization}, 2nd ed., 2006, ch.~6). In
practice one updates the inverse $B_k^{-1}$ directly. L-BFGS (Liu and Nocedal, Mathematical
Programming, 1989) keeps only the last $m\approx10$ pairs $(s_k, y_k)$ and applies $B_k^{-1}$ implicitly in $O(mp)$
operations. These methods converge superlinearly near a well-conditioned minimum with
exact gradients; with mini-batch noise the secant pairs become unreliable, which is why
they are used for full-batch problems of moderate size (\texttt{torch.optim.LBFGS},
\texttt{scipy.optimize.minimize(method="L-BFGS-B")}) rather than for training large
networks.

\begin{chaptersummary}
\begin{itemize}[nosep]
  \item On a quadratic, gradient descent needs $\eta<2/L$ and converges at rate $\frac{\kappa - 1}{\kappa + 1}$ at best: about $\kappa/2$
        steps per factor $e$. At a critical point a definite Hessian decides min or max; mixed
        signs give a saddle.
  \item Heavy-ball momentum, tuned, achieves $\frac{\sqrt\kappa - 1}{\sqrt\kappa + 1}$ on quadratics; Nesterov's method achieves
        $1 - 1/\sqrt\kappa$ for all smooth strongly convex functions.
  \item AdaGrad, RMSProp and Adam scale each coordinate by the root of an average of squared
        gradients; Adam adds momentum and a bias correction $1/(1 - \beta^k)$; its first step is $-\eta\sign(g)$.
        AdamW decouples weight decay from the scaling.
  \item Schedules decrease $\eta$ (step, exponential, cosine), often after a warm-up.
  \item BFGS and L-BFGS approximate Newton's method from gradient differences; they need
        $y^\top s>0$ and accurate gradients.
\end{itemize}
\end{chaptersummary}

\section*{Exercises}
\addcontentsline{toc}{section}{Exercises}

\begin{exercise}\label{optimizers:exr:1}
For $\ell = \frac12(\theta_1^2 + 25\theta_2^2)$, find the largest learning rate for which gradient descent converges, and
describe the iterates for $\eta = 0.08$ and for $\eta = 1/25$.
\solutionlink{optimizers:exr:1}
\end{exercise}

\begin{exercise}\label{optimizers:exr:2}
Show that the PyTorch and classical forms of heavy-ball momentum produce the same iterates
when $\eta$ is constant, and both equal Polyak's form. What changes if $\eta$ is divided by $10$ at some
step?
\solutionlink{optimizers:exr:2}
\end{exercise}

\begin{exercise}\label{optimizers:exr:3}
Show that the first Adam step is $-\eta g_1/(\abs{g_1} + \epsilon)$, and that Adam's update is unchanged (for $\epsilon = 0$) if
every gradient is replaced by $Dg_k$ for a fixed diagonal $D$ with positive entries.
\solutionlink{optimizers:exr:3}
\end{exercise}

\begin{exercise}\label{optimizers:exr:4}
With $\beta_2 = 0.999$, what is $\E[v_1]$ in terms of $\E[g^2]$ without bias correction? After how many steps is the
bias factor $1 - \beta_2^k$ above $0.5$? Above $0.99$?
\solutionlink{optimizers:exr:4}
\end{exercise}

\begin{exercise}\label{optimizers:exr:5}
Check that the BFGS matrix $B_{k+1}$ satisfies the secant equation $B_{k+1}s_k = y_k$, and that
$v^\top B_{k+1}v>0$ for $v = s_k$ when $B_k$ is positive definite and $y_k^\top s_k>0$.
\solutionlink{optimizers:exr:5}
\end{exercise}
\end{notechapter}
